\documentclass[UTF8,a4paper,11pt]{ctexart}

\usepackage{amsmath,amssymb,bm,mathtools}
\usepackage{booktabs,longtable}
\usepackage{enumitem}
\usepackage{geometry}
\usepackage{xcolor}
\usepackage{hyperref}

\geometry{left=24mm,right=24mm,top=24mm,bottom=26mm}
\hypersetup{colorlinks=true,linkcolor=blue!55!black,urlcolor=blue!55!black}
\setlist{nosep}
\newcommand{\R}{\mathbb{R}}
\newcommand{\SO}{\mathrm{SO}(3)}
\newcommand{\Exp}{\operatorname{Exp}}
\newcommand{\Log}{\operatorname{Log}}
\newcommand{\diag}{\operatorname{diag}}
\newcommand{\clip}{\operatorname{clip}}
% Do not call this command \skew: LaTeX already defines \skew for adjusting
% math accents, and overriding it breaks amsmath's align environment.
\newcommand{\skewmat}[1]{\left\lfloor #1\right\rfloor_\times}
\newcommand{\normQ}[2]{\left\lVert #1\right\rVert_{#2}^{2}}
\newcommand{\pseudocode}[1]{%
  \begin{center}
  \fbox{\begin{minipage}{0.94\linewidth}\small #1\end{minipage}}
  \end{center}}

\title{GPT 四旋翼流形 MPC 与轨迹优化\\数学原理及实现伪代码}
\author{对应工程文件：\texttt{gptmpc.cpp/.h} 与 \texttt{gpttraj.cpp/.h}}
\date{2026年8月}

\begin{document}
\maketitle
\tableofcontents

\section{文档范围}

本文说明工程中以下四个文件所实现的算法，而不是对参考论文的逐字复述：
\begin{itemize}
  \item \texttt{src/realflight\_modules/px4ctrl/include/px4ctrl/gptmpc.h}；
  \item \texttt{src/realflight\_modules/px4ctrl/src/gptmpc.cpp}；
  \item \texttt{src/realflight\_modules/px4ctrl/include/px4ctrl/gpttraj.h}；
  \item \texttt{src/realflight\_modules/px4ctrl/src/gpttraj.cpp}。
\end{itemize}

控制器以总推力加速度和机体角速度为直接输入，在
$\R^3\times\R^3\times\SO$ 上构造误差状态 MPC。轨迹优化器使用逐次凸化
（Sequential Convex Programming, SCP），每轮以 Eigen 构造稀疏二次规划，
再由 OSQP 求解。

\section{坐标系、状态与输入}

\subsection{工程坐标系}

工程内部采用 ENU 惯性坐标系和 FLU 机体系。姿态矩阵
\begin{equation}
  \bm R\in\SO
\end{equation}
将机体系向量主动旋转到惯性系；其第三列
$\bm R\bm e_3$ 是正总推力在惯性系中的方向，其中
$\bm e_3=[0,0,1]^\mathsf T$。这与 Lu.pdf 中采用的 FRD 机体系不同，
因此本实现中的推力项及姿态雅可比符号也相应改变。

\subsection{状态和控制量}

无人机状态和控制输入定义为
\begin{align}
  \bm x &= (\bm p,\bm v,\bm R)
  \in \mathcal M=\R^3\times\R^3\times\SO, \\
  \bm u &=
  \begin{bmatrix}a_T&\bm\omega^\mathsf T\end{bmatrix}^{\mathsf T}
  \in\R^4,
\end{align}
其中 $\bm p$、$\bm v$ 是惯性系位置和速度，$a_T$ 是沿机体 $+z$
方向的总推力加速度，$\bm\omega$ 是 FLU 机体系角速度。

忽略气动阻力时，连续动力学为
\begin{align}
  \dot{\bm p} &= \bm v, \\
  \dot{\bm v} &= -g\bm e_3+a_T\bm R\bm e_3, \\
  \dot{\bm R} &= \bm R\skewmat{\bm\omega}.
  \label{eq:continuous_model}
\end{align}
这里 $g>0$。控制器最终输出给工程底层的是牛顿制总推力
\begin{equation}
  T=m a_T,
  \label{eq:thrust_conversion}
\end{equation}
以及角速度 $\bm\omega$。

\subsection{离散传播}

采样周期为 $\Delta t$ 时，实际代码采用
\begin{align}
  \bm p_{k+1} &= \bm p_k+\Delta t\,\bm v_k,\\
  \bm v_{k+1} &= \bm v_k+\Delta t
  \left(-g\bm e_3+a_{T,k}\bm R_k\bm e_3\right),\\
  \bm R_{k+1} &= \bm R_k\Exp(\Delta t\,\bm\omega_k).
  \label{eq:discrete_model}
\end{align}
最后一式通过指数映射传播旋转，因此不会破坏
$\bm R^\mathsf T\bm R=\bm I$ 和 $\det\bm R=1$。

\section{$\SO$ 上的局部坐标}

\subsection{帽算子与指数/对数映射}

对 $\bm\phi=[\phi_x,\phi_y,\phi_z]^\mathsf T$，定义
\begin{equation}
  \skewmat{\bm\phi}=
  \begin{bmatrix}
    0&-\phi_z&\phi_y\\
    \phi_z&0&-\phi_x\\
    -\phi_y&\phi_x&0
  \end{bmatrix}.
\end{equation}
旋转向量通过 Rodrigues 公式映射为旋转矩阵：
\begin{equation}
  \Exp(\bm\phi)=
  \bm I+\frac{\sin\theta}{\theta}\skewmat{\bm\phi}
  +\frac{1-\cos\theta}{\theta^2}\skewmat{\bm\phi}^{2},
  \qquad \theta=\lVert\bm\phi\rVert_2.
\end{equation}

\subsection{流形加减法}

对九维局部增量
\begin{equation}
  \delta\bm x=
  \begin{bmatrix}
    \delta\bm p^\mathsf T&
    \delta\bm v^\mathsf T&
    \delta\bm\theta^\mathsf T
  \end{bmatrix}^{\mathsf T}\in\R^9,
\end{equation}
本实现使用右扰动：
\begin{equation}
  \bm x\boxplus\delta\bm x=
  \left(
    \bm p+\delta\bm p,
    \bm v+\delta\bm v,
    \bm R\Exp(\delta\bm\theta)
  \right).
  \label{eq:boxplus}
\end{equation}
两个状态之间的最小局部误差为
\begin{equation}
  \bm x\boxminus\bm x^d=
  \begin{bmatrix}
    \bm p-\bm p^d\\
    \bm v-\bm v^d\\
    \Log\!\left((\bm R^d)^\mathsf T\bm R\right)
  \end{bmatrix}\in\R^9.
  \label{eq:boxminus}
\end{equation}
因此优化变量只有九维，不需要四元数单位模约束，也不存在欧拉角奇异点。

\section{流形模型预测控制}

\subsection{参考轨迹}

控制器接受三种参考来源，优先级如下：
\begin{enumerate}
  \item 由 \texttt{setTrajectory()} 设置的优化轨迹；
  \item 由 \texttt{setReferencePreview()} 设置的至少 $N+1$ 个参考点；
  \item 当前控制周期传入的单个 \texttt{Ref\_State\_t}。
\end{enumerate}

对第三种情况，代码用常加速度外推
\begin{align}
  \bm p_k^d &= \bm p^d+t_k\bm v^d+\frac{1}{2}t_k^2\bm a^d,\\
  \bm v_k^d &= \bm v^d+t_k\bm a^d,
\end{align}
并令期望比力为
\begin{equation}
  \bm f_k^d=\bm a^d+t_k\bm j^d+g\bm e_3,
  \qquad a_{T,k}^d=\clip
  \left(\lVert\bm f_k^d\rVert_2,a_{T,\min},a_{T,\max}\right).
\end{equation}
以 $\bm z_B^d=\bm f_k^d/\lVert\bm f_k^d\rVert_2$ 和参考偏航角构造
$\bm R_k^d$。参考角速度由相邻参考姿态差分得到：
\begin{equation}
  \bm\omega_k^d=\frac{1}{\Delta t}
  \Log\!\left((\bm R_k^d)^\mathsf T\bm R_{k+1}^d\right).
\end{equation}

\subsection{仿射误差动力学}

记离散非线性传播映射为
\begin{equation}
  \bm x_{k+1}=F_d(\bm x_k,\bm u_k).
\end{equation}
在参考点 $(\bm x_k^d,\bm u_k^d)$ 附近定义
\begin{align}
  \delta\bm x_k &= \bm x_k\boxminus\bm x_k^d,\\
  \delta\bm u_k &= \bm u_k-\bm u_k^d.
\end{align}
则下一时刻误差映射为
\begin{equation}
  \mathcal E_k(\delta\bm x_k,\delta\bm u_k)=
  F_d(\bm x_k^d\boxplus\delta\bm x_k,
      \bm u_k^d+\delta\bm u_k)
  \boxminus\bm x_{k+1}^d.
\end{equation}

代码使用前向有限差分计算雅可比：
\begin{align}
  \bm A_k(:,i)&\approx
  \frac{\mathcal E_k(\epsilon\bm e_i,\bm0)
        -\mathcal E_k(\bm0,\bm0)}{\epsilon},\\
  \bm B_k(:,j)&\approx
  \frac{\mathcal E_k(\bm0,\epsilon\bm e_j)
        -\mathcal E_k(\bm0,\bm0)}{\epsilon},\\
  \bm c_k&=\mathcal E_k(\bm0,\bm0),
  \qquad \epsilon=10^{-6}.
\end{align}
从而得到
\begin{equation}
  \delta\bm x_{k+1}
  =\bm A_k\delta\bm x_k
  +\bm B_k\delta\bm u_k+\bm c_k.
  \label{eq:affine_error_model}
\end{equation}
$\bm c_k$ 是参考轨迹的动力学缺陷；即使输入的参考并非严格动力学可行，
它仍会进入预测模型，而不会被静默忽略。

\subsection{预测模型凝聚}

递推式~\eqref{eq:affine_error_model} 可写成
\begin{equation}
  \delta\bm x_k=\bm\Phi_k\delta\bm x_0
  +\sum_{j=0}^{k-1}\bm\Gamma_{k,j}\delta\bm u_j
  +\bm\eta_k.
\end{equation}
堆叠未来状态和输入：
\begin{align}
  \delta\bm X&=
  \begin{bmatrix}
    \delta\bm x_1^\mathsf T&\cdots&\delta\bm x_N^\mathsf T
  \end{bmatrix}^{\mathsf T},\\
  \delta\bm U&=
  \begin{bmatrix}
    \delta\bm u_0^\mathsf T&\cdots&\delta\bm u_{N-1}^\mathsf T
  \end{bmatrix}^{\mathsf T},
\end{align}
得到凝聚预测模型
\begin{equation}
  \delta\bm X
  =\bm S_x\delta\bm x_0+\bm S_u\delta\bm U+\bm s_c.
  \label{eq:condensed_prediction}
\end{equation}

\subsection{代价函数和 QP}

有限时域代价为
\begin{equation}
  J=\sum_{k=1}^{N}\normQ{\delta\bm x_k}{\bm Q_k}
  +\sum_{k=0}^{N-1}\normQ{\delta\bm u_k}{\bm R_k}.
\end{equation}
其中默认 $N=8$，状态权重和输入权重取
\begin{align}
  \bm Q_k&=\diag(15000,15000,15000,40,40,40,80,80,80),\\
  \bm R_k&=\diag(0.5,0.6,0.6,0.6).
\end{align}

令 $\bm z=\bm S_x\delta\bm x_0+\bm s_c$，代入
式~\eqref{eq:condensed_prediction}，可化为 OSQP 使用的标准形式
\begin{equation}
  \min_{\delta\bm U}\quad
  \frac12\delta\bm U^\mathsf T\bm H\delta\bm U
  +\bm q^\mathsf T\delta\bm U,
\end{equation}
其中
\begin{align}
  \bm H&=2(\bm S_u^\mathsf T\bar{\bm Q}\bm S_u+\bar{\bm R}),\\
  \bm q&=2\bm S_u^\mathsf T\bar{\bm Q}\bm z.
\end{align}

控制约束转换为增量约束：
\begin{align}
  a_{T,\min}-a_{T,k}^d
  &\leq \delta a_{T,k}
  \leq a_{T,\max}-a_{T,k}^d,\\
  -\bm\omega_{\max}-\bm\omega_k^d
  &\leq \delta\bm\omega_k
  \leq \bm\omega_{\max}-\bm\omega_k^d.
\end{align}
求解后只执行第一项，即滚动时域策略
\begin{equation}
  \bm u_{\mathrm{cmd}}=\bm u_0^d+\delta\bm u_0^\star.
\end{equation}

\subsection{MPC 伪代码}

\pseudocode{
\textbf{算法 1：\texttt{GptMpcControl::calculateControl}}\\[2mm]
\textbf{输入：}当前状态 $(\bm p,\bm v,\bm R)$，参考状态/轨迹，参数 $m,g$，周期 $\Delta t$\\
\textbf{输出：}总推力 $T$，机体角速度 $\bm\omega_{\rm cmd}$
\begin{enumerate}[label=\arabic*.]
  \item 若处于手动模式：计算姿态误差比例角速度，按油门映射总推力，然后返回。
  \item 从优化轨迹、预览序列或单点参考中构造 $N+1$ 个参考状态及 $N$ 个参考输入。
  \item 用式~\eqref{eq:boxminus} 计算当前九维误差 $\delta\bm x_0$。
  \item 对 $k=0,\ldots,N-1$：
  \begin{enumerate}[label*=\arabic*.]
    \item 用式~\eqref{eq:discrete_model} 定义流形误差传播 $\mathcal E_k$；
    \item 以有限差分计算 $\bm A_k,\bm B_k$；
    \item 保存参考动力学缺陷 $\bm c_k$。
  \end{enumerate}
  \item 递推构造 $\bm S_x,\bm S_u,\bm s_c$。
  \item 构造 Hessian $\bm H$、梯度 $\bm q$ 以及输入上下界。
  \item 用 Eigen 转为 CSC 稀疏矩阵，调用 OSQP 求解 $\delta\bm U^\star$。
  \item 若求解成功，取 $\bm u_0^d+\delta\bm u_0^\star$；否则保持上一周期输入并重新限幅。
  \item 按式~\eqref{eq:thrust_conversion} 输出 $T=m a_T$ 和 $\bm\omega_{\rm cmd}$。
\end{enumerate}}

\section{SCP 轨迹优化}

\subsection{优化目标}

给定初始状态、终端状态和若干中间航点，轨迹优化器寻找满足
式~\eqref{eq:discrete_model} 及输入边界的短时间轨迹。当前实现将总时间搜索
与固定时间下的 SCP 分开：外层逐步缩短 $t_f$，内层对每个给定的 $t_f$
求解一系列凸 QP。

时间均匀离散为 $n$ 段：
\begin{equation}
  \Delta t=\frac{t_f}{n},\qquad k=0,\ldots,n.
\end{equation}

\subsection{初始轨迹}

令折线路径依次连接起点、中间航点和终点，总弧长为 $L$。初始位置按弧长
均匀插值；内部节点速度用中心差分：
\begin{equation}
  \bm v_k=\frac{\bm p_{k+1}-\bm p_{k-1}}{2\Delta t}.
\end{equation}
再由速度差分估计加速度，并用
$\bm a_k+g\bm e_3$ 构造推力方向和初始姿态。初始控制量为
\begin{align}
  a_{T,k}&=\clip\left(
  \left(\frac{\bm v_{k+1}-\bm v_k}{\Delta t}+g\bm e_3\right)^\mathsf T
  \bm R_k\bm e_3,
  a_{T,\min},a_{T,\max}\right),\\
  \bm\omega_k&=\frac1{\Delta t}
  \Log(\bm R_k^\mathsf T\bm R_{k+1}).
\end{align}

\subsection{每轮 SCP 子问题}

第 $\ell$ 轮的名义轨迹为
$(\bm x_k^\ell,\bm u_k^\ell)$。优化变量为
\begin{equation}
  \bm y=\left\{
  \delta\bm x_0,\ldots,\delta\bm x_n,
  \delta\bm u_0,\ldots,\delta\bm u_{n-1},
  \bm d_0,\ldots,\bm d_{n-1}
  \right\},
\end{equation}
其中 $\bm d_k\in\R^9$ 是虚拟控制，用于防止线性化子问题人为不可行。

在名义轨迹处用与 MPC 相同的方法计算 $\bm A_k,\bm B_k,\bm c_k$，
并施加
\begin{equation}
  \delta\bm x_{k+1}
  =\bm A_k\delta\bm x_k
  +\bm B_k\delta\bm u_k+\bm c_k+\bm d_k.
  \label{eq:scp_dynamics}
\end{equation}

边界条件为
\begin{align}
  \delta\bm x_0&=\bm0,\\
  \delta\bm x_n&=\bm x_f\boxminus\bm x_n^\ell.
\end{align}

当前代码按折线累计弧长为中间航点分配有序离散节点 $\kappa_j$：
\begin{equation}
  \kappa_j\approx
  \operatorname{round}\left(n\frac{L_j}{L}\right),
  \qquad 1\leq\kappa_1<\cdots<\kappa_M\leq n-1.
\end{equation}
球形容差在 QP 中用保守的轴对齐盒约束实现：
\begin{equation}
  -\frac{r_j}{\sqrt3}\bm1
  \leq
  \bm p_{\kappa_j}^\ell+\delta\bm p_{\kappa_j}-\bm p_j^w
  \leq
  \frac{r_j}{\sqrt3}\bm1.
  \label{eq:waypoint_box}
\end{equation}
若式~\eqref{eq:waypoint_box} 成立，则必有
$\lVert\bm p_{\kappa_j}-\bm p_j^w\rVert_2\leq r_j$。

信赖域逐分量表示为
\begin{align}
  \lVert\delta\bm p_k\rVert_\infty&\leq\rho_p,\\
  \lVert\delta\bm v_k\rVert_\infty&\leq\rho_v,\\
  \lVert\delta\bm\theta_k\rVert_\infty&\leq\rho_R.
\end{align}
输入增量边界与 MPC 相同。

实际实现的凸二次代价为
\begin{align}
  J_{\rm SCP}={}&
  w_x\sum_{k=0}^{n}\lVert\delta\bm x_k\rVert_2^2
  +w_u\sum_{k=0}^{n-1}\lVert\delta\bm u_k\rVert_2^2\\
  &+w_{\Delta u}\sum_{k=1}^{n-1}
  \lVert\delta\bm u_k-\delta\bm u_{k-1}\rVert_2^2
  +w_d\sum_{k=0}^{n-1}\lVert\bm d_k\rVert_2^2.
  \label{eq:scp_cost}
\end{align}
$w_d$ 取很大值，使虚拟控制仅在必要时出现。QP 解通过
\begin{align}
  \bm p_k^{\ell+1}&=\bm p_k^\ell+\delta\bm p_k^\star,\\
  \bm v_k^{\ell+1}&=\bm v_k^\ell+\delta\bm v_k^\star,\\
  \bm R_k^{\ell+1}&=\bm R_k^\ell\Exp(\delta\bm\theta_k^\star),\\
  \bm u_k^{\ell+1}&=\bm u_k^\ell+\delta\bm u_k^\star
\end{align}
映射回原状态空间。

收敛条件同时检查最大状态更新和虚拟控制：
\begin{equation}
  \max_k\lVert\delta\bm x_k^\star\rVert_\infty<\epsilon_x,
  \qquad
  \max_k\lVert\bm d_k^\star\rVert_\infty<\epsilon_d.
\end{equation}

\subsection{飞行时间搜索}

初始总时间由路径长度估计：
\begin{equation}
  t_f^{(0)}=\clip
  \left(1.5\frac{L}{v_{\rm init}},t_{\min},t_{\max}\right).
\end{equation}
若该时间不可行，则逐步放大；找到第一个可行时间后，以缩放因子
$\gamma\in(0,1)$ 尝试
\begin{equation}
  t_{\rm trial}=\max(t_{\min},\gamma t_{\rm feasible}),
\end{equation}
并结合可行/不可行时间区间的中点更新搜索。这里的 ``可行'' 指 OSQP 成功，
且 SCP 结束时最大虚拟控制不超过给定阈值。

\subsection{轨迹优化伪代码}

\pseudocode{
\textbf{算法 2：固定总时间下的 SCP}\\[2mm]
\textbf{输入：}边界状态、航点、总时间 $t_f$、离散段数 $n$\\
\textbf{输出：}动力学可行轨迹或失败状态
\begin{enumerate}[label=\arabic*.]
  \item 沿起点--航点--终点折线初始化 $\bm p_k,\bm v_k,\bm R_k,\bm u_k$。
  \item 依据累计弧长计算严格递增的航点节点 $\kappa_j$。
  \item 对 $\ell=0,\ldots,L_{\rm SCP}-1$：
  \begin{enumerate}[label*=\arabic*.]
    \item 沿名义轨迹有限差分计算 $\bm A_k,\bm B_k,\bm c_k$；
    \item 创建变量 $\delta\bm x_k,\delta\bm u_k,\bm d_k$；
    \item 按式~\eqref{eq:scp_cost} 构造稀疏 Hessian；
    \item 加入边界、式~\eqref{eq:scp_dynamics}、航点、信赖域及输入边界约束；
    \item 转换为 $\bm l\leq\bm A\bm y\leq\bm u$，调用 OSQP；
    \item 若 OSQP 失败，返回不可行；
    \item 用 $\boxplus$ 更新名义状态，用向量加法更新输入；
    \item 若状态更新及虚拟控制均低于阈值，则终止 SCP。
  \end{enumerate}
  \item 若最大虚拟控制小于动力学容差，返回成功；否则返回不可行。
\end{enumerate}}

\pseudocode{
\textbf{算法 3：外层时间搜索}\\[2mm]
\begin{enumerate}[label=\arabic*.]
  \item 根据路径长度和初始速度估计 $t_f^{(0)}$。
  \item 调用算法 2；若失败且未达到 $t_{\max}$，放大 $t_f$ 后重试。
  \item 保存当前可行轨迹和可行时间上界。
  \item 在规定次数内缩短飞行时间，并再次调用算法 2：
  \begin{itemize}
    \item 成功：保存更短轨迹，收紧可行时间上界；
    \item 失败：提高不可行时间下界。
  \end{itemize}
  \item 返回找到的最短可行轨迹。
\end{enumerate}}

\section{OSQP 与 Eigen 数据接口}

两种算法最终都求解
\begin{equation}
  \begin{aligned}
  \min_{\bm y}\quad &\frac12\bm y^\mathsf T\bm P\bm y+\bm q^\mathsf T\bm y,\\
  \text{s.t.}\quad &\bm l\leq\bm A\bm y\leq\bm u.
  \end{aligned}
\end{equation}
实现步骤为：
\begin{enumerate}
  \item 使用 Eigen 稠密或 \texttt{Triplet} 形式组装矩阵；
  \item 转换为列压缩 CSC 存储；
  \item $\bm P$ 只向 OSQP 传入上三角部分；
  \item 设置缩放、warm start、polishing 和收敛容差；
  \item 接受 \texttt{solved} 或 \texttt{solved inaccurate} 状态；
  \item 将 OSQP 解复制回 Eigen 向量并更新控制或轨迹。
\end{enumerate}

\section{实现边界与论文完整框架的区别}

为避免把尚未实现的论文内容误写成现有功能，当前代码边界明确如下：
\begin{itemize}
  \item 已实现流形误差 MPC、推力/角速度直接输入、SCP 轨迹优化、输入边界、
  信赖域、虚拟控制、外层时间搜索和逐轮轨迹历史记录。
  \item 已把 $\lambda_k^j$、$\mu_k^j$、空间 CSTC 与严格顺序 CSTC 加入 QP；
  基于折线弧长的节点仅用于 one-hot 初始化，航点经过时刻由优化器选择。
  \item 轨迹成功条件同时检查动力学缺陷、航点互补残差、顺序互补残差和
  CSTC 松弛量，OSQP 单轮成功本身不被当作原非凸问题成功。
  \item 当前轨迹子问题使用虚拟控制的二范数平方惩罚，而论文公式给出的是
  一范数惩罚；二范数形式可直接形成纯 QP。
  \item 当前输入约束是总推力加速度和机体角速度盒约束，没有直接施加单电机
  转速、力矩或旋翼分配约束。
  \item 四个 GPT 文件未实现 CBF/HOCBF 避障、气动阻力补偿和 INDI 内环；
  工程已有的角速度底层仍负责后续力矩控制及电机分配。
  \item 有限差分雅可比提高了实现通用性，但比手工解析雅可比增加计算量。
\end{itemize}

\section{建议的调用流程}

\begin{enumerate}
  \item 创建 \texttt{GptTrajectoryOptimizer}，设置离散段数、推力和角速度边界。
  \item 调用 \texttt{optimize(initial, terminal, waypoints)} 得到
  \texttt{GptTrajectoryResult}。
  \item 检查 \texttt{success}、动力学/航点/顺序残差、CSTC 松弛量和总时间。
  \item 调用控制器的
  \texttt{setTrajectory(result, start\_time)} 设置参考轨迹。
  \item 每个外环周期调用与现有控制器兼容的 \texttt{calculateControl()}。
  \item 将输出的牛顿制总推力和 FLU 机体系角速度发送给现有底层角速度节点。
\end{enumerate}

\end{document}
